% ================= MỤC 1 =================
\section{Ứng dụng}

\begin{frame}{Phép tích chập là gì?}
    Phép tích chập (Convolution) đo lường mức độ chồng lấp giữa hai hàm số khi một hàm được lật ngược và dịch chuyển qua hàm kia.
    \vspace{0.5cm}
    
    \textbf{Tích chập tuyến tính} cho hai dãy rời rạc $x[n]$ và $h[n]$:
    $$(x * h)[n] = \sum_{m=-\infty}^{\infty} x[m] h[n - m]$$
    \vspace{0.3cm}
    
    \textbf{Tích chập vòng} (áp dụng cho dãy độ dài hữu hạn $N$):
    $$(x \circledast h)[n] = \sum_{m=0}^{N-1} x[m] h[(n - m) \bmod N]$$
\end{frame}

\begin{frame}{Định lý Tích chập: Chìa khóa chuyển miền}
    Tính toán trực tiếp tích chập cần độ phức tạp $\mathcal{O}(N^2)$. Tuy nhiên, Định lý Tích chập tạo ra bước ngoặt lớn:
    \vspace{0.5cm}
    
    \begin{block}{Định lý Tích chập}
        Phép tích chập vòng trong miền thời gian/không gian hoàn toàn tương đương với phép \textbf{nhân từng phần tử} trong miền tần số.
        $$\text{DFT}\{x \circledast h\} = \text{DFT}\{x\} \cdot \text{DFT}\{h\}$$
    \end{block}
    \vspace{0.5cm}
    
    Thay vì tính $\mathcal{O}(N^2)$ ở không gian gốc, ta tính $\mathcal{O}(N)$ ở không gian tần số. Vấn đề còn lại là: \textit{Làm sao chuyển miền đủ nhanh?}
\end{frame}

\begin{frame}{Thuật toán Cooley-Tukey (FFT)}
    \begin{itemize}
        \item \textbf{Ý tưởng:} Chia để trị (Divide and Conquer).
        \item Chia bài toán DFT kích thước $N$ thành các bài toán chẵn - lẻ kích thước $N/2$.
        \item Khai thác tính đối xứng của nghiệm phức đơn vị: $\omega_N = e^{-j2\pi/N}$.
    \end{itemize}
    \vspace{0.5cm}
    
    \begin{alertblock}{Kết quả}
        Giảm độ phức tạp chuyển miền (DFT) từ $\mathcal{O}(N^2)$ xuống chỉ còn $\mathcal{O}(N \log N)$.
    \end{alertblock}
\end{frame}

\begin{frame}{Quy trình Tích chập nhanh }
    Để tính $(x * h)$ bằng FFT, ta thực hiện 4 bước sau:
    \vspace{0.3cm}
    \begin{enumerate}
        \item \textbf{Đệm 0 (Zero-padding):} Đệm thêm các số $0$ cho $x$ và $h$ để đạt chiều dài $N$ (thường là lũy thừa của 2).
        \item \textbf{Chuyển miền (FFT):} Tính $X = \text{FFT}(x)$ và $H = \text{FFT}(h)$. \\ $\Rightarrow$ Tốn $\mathcal{O}(N \log N)$.
        \item \textbf{Nhân phổ:} Lấy $Y[k] = X[k] \cdot H[k]$ với mọi $k$. \\ $\Rightarrow$ Tốn $\mathcal{O}(N)$.
        \item \textbf{Dịch ngược (IFFT):} Khôi phục kết quả bằng $\text{IFFT}(Y)$. \\ $\Rightarrow$ Tốn $\mathcal{O}(N \log N)$.
    \end{enumerate}
    \vspace{0.3cm}
    \textbf{Tổng thời gian:} $\mathcal{O}(N \log N)$ thay vì $\mathcal{O}(N^2)$.
\end{frame}

% ================= MỤC 2 =================
\begin{frame}{Nhân đa thức: Bản chất là Tích chập}
    Cho hai đa thức $A(x) = \sum_{i=0}^{n-1} a_i x^i$ và $B(x) = \sum_{j=0}^{m-1} b_j x^j$.
    \vspace{0.3cm}
    
    Đa thức tích $C(x) = A(x)B(x)$ có hệ số $c_k$ được tính bằng:
    $$c_k = \sum a_i b_{k-i}$$
    
    \begin{itemize}
        \item Biểu thức trên \textbf{chính xác} là định nghĩa của phép tích chập tuyến tính giữa hai dãy hệ số: $(a * b)$.
        \item \textbf{Giải pháp:} Đệm $0$ cho hai dãy đạt chiều dài $N \ge n+m-1$, rồi áp dụng quy trình \textit{Fast Convolution} bằng thuật toán Cooley-Tukey.
        \item \textbf{Ứng dụng:} Mật mã học (RSA), tính toán số nguyên cực lớn, đại số máy tính.
    \end{itemize}
\end{frame}

\begin{frame}{Xử lý ảnh: Phép lọc và Tích chập 2D}
    Trong ảnh số, làm mượt (blur), tìm cạnh (edge detection) đều là \textbf{phép tích chập 2 chiều} giữa:
    \begin{itemize}
        \item Ma trận điểm ảnh ban đầu.
        \item Ma trận bộ lọc (Filter kernel / Mask).
    \end{itemize}
    \vspace{0.5cm}
    
    \textbf{Vấn đề:} Trượt một bộ lọc lớn qua từng pixel tiêu tốn tài nguyên khổng lồ.
    \vspace{0.3cm}
    
    \textbf{Giải pháp bằng FFT 2D:}
    \begin{enumerate}
        \item Biến đổi cả ảnh và bộ lọc sang phổ tần số không gian (dùng FFT 1D trên hàng, sau đó trên cột).
        \item Nhân từng phần tử hai ma trận tần số.
        \item Dùng IFFT 2D khôi phục ảnh.
    \end{enumerate}
\end{frame}

\begin{frame}{Xử lý âm thanh: Lọc và Hiệu ứng}
    Các hiệu ứng (reverb, echo) hay lọc tạp âm (lọc thông thấp, thông cao) chính là \textbf{phép tích chập 1 chiều} giữa:
    \begin{itemize}
        \item Tín hiệu âm thanh đầu vào.
        \item Đáp ứng xung (Impulse response) của bộ lọc.
    \end{itemize}
    \vspace{0.5cm}
    
    \textbf{Tối ưu cho thời gian thực (Real-time):}
    \begin{itemize}
        \item Việc tính tích chập trực tiếp trên đoạn âm thanh dài là bất khả thi.
        \item \textbf{Kỹ thuật:} Cắt âm thanh thành các khung (frames).
        \item Kết hợp các phương pháp \textit{Overlap-Add} hoặc \textit{Overlap-Save} với FFT Cooley-Tukey.
        \item \textbf{Kết quả:} Lọc nhiễu, xử lý giọng nói tức thời (Zoom, Google Meet, trợ lý ảo) mà không bị trễ.
    \end{itemize}
\end{frame}